\section{Tools and Miscellaneous}

The modern technical approach to cybersecurity marks a fundamental shift from a static "Security" model to a dynamic \textbf{Resilience} model. In this framework, security is no longer defined by the absence of incidents, but by the organization's capacity to maintain functions despite ongoing attacks.\\
For this reason, it become essential the defense in depth paradigm.

\subsubsection*{Defense in Depth}
Operational security is structured into multiple, redundant layers. This strategy ensures that the failure of a single control point (e.g., a firewall) does not lead to a total system compromise.
\begin{itemize}
    \item \textbf{Layered Structure:} The protection extends from the \textbf{Data} core outwards through \textbf{Applications}, \textbf{Hosts}, \textbf{Network}, and finally the \textbf{Physical/Perimeter} layer.
    \item \textbf{The Human Factor:} The user is identified as the outermost and often weakest layer, requiring a balance between technical controls and awareness training.
\end{itemize}

\subsection{Operational Monitoring and Log Management}
Without comprehensive monitoring, it is impossible to detect anomalies or reconstruct the timeline of an attack.
\begin{itemize}
    \item \textbf{Log Analysis:} Logs serve both as a tool for \textbf{Incident Response} and as a \textbf{Compliance} requirement (e.g., GDPR, NIS2). They must be centralized and protected against tampering.
    \item \textbf{SIEM Systems:} Security Information and Event Management (SIEM) tools are utilized to correlate data across different layers, identifying patterns that individual logs might miss.
\end{itemize}

\subsection{Business Continuity and Disaster Recovery}
Even if a system is monitored and protected, it is inevitable that some attacks will succeed. The focus then shifts to minimizing the impact and restoring operations as quickly as possible.
\begin{itemize}
    \item \textbf{Business Continuity (BC):} The strategic capability of an organization to continue delivery of products or services at acceptable predefined levels following a disruptive incident.
    \item \textbf{Disaster Recovery (DR):} The specific technical processes and infrastructures (e.g., secondary data centers, backup restores) required to regain access to IT systems.
    \item \textbf{RPO vs. RTO:} 
    \begin{itemize}
        \item \textbf{RPO (Recovery Point Objective):} Maximum tolerable data loss (measured in time).
        \item \textbf{RTO (Recovery Time Objective):} Maximum tolerable downtime before the service must be restored.
    \end{itemize}
\end{itemize}

\begin{figure}[H]
\centering
\begin{tikzpicture}[font=\small,node distance=6mm,>=Latex,
box/.style={rectangle,rounded corners,draw=black,fill=gray!10,align=center,text width=5cm,minimum height=6mm}]
\node[box] (event) {\textbf{Disruptive Event}\\System Failure / Attack};
\node[box, below=of event] (dr) {\textbf{Disaster Recovery}\\Technical Restoration};
\node[box, below=of dr, fill=blue!5] (bc) {\textbf{Business Continuity}\\Service Resumption};

\draw[->] (event) -- (dr);
\draw[->] (dr) -- node[right]{Restoration of Assets} (bc);
\end{tikzpicture}
\caption{The relationship between technical recovery (DR) and business stability (BC).}
\end{figure}

\noindent
A proactive stance requires the continuous identification of threats.
\begin{itemize}
    \item \textbf{Vulnerability Scans:} Automated tools used to identify "open windows" in the infrastructure before attackers can exploit them.
    \item \textbf{Zero-Day Awareness:} The acknowledgment that some vulnerabilities are unknown to vendors, necessitating a focus on \textit{behavioral} monitoring rather than just signature-based detection.
\end{itemize}



\subsection{Methodology for Information Security Risk Assessment}
The technical framework for risk management follows a structured, iterative process designed to align organizational assets with the threat landscape. 

\subsubsection*{1. Identification of Assets and Vulnerabilities}
The first phase involves mapping the surface area of the organization:
\begin{itemize}
    \item \textbf{Asset Inventory:} Differentiation between \textbf{Primary Assets} (core business processes and information) and \textbf{Supporting Assets} (infrastructure, software, personnel, sites).
    \item \textbf{Vulnerability Mapping:} Identifying technical or organizational weaknesses (e.g., unpatched software, lack of physical access control) that could be exploited by a threat.
\end{itemize}

\subsubsection*{2. Risk Evaluation: Impact and Likelihood}
Risk is not a static value but a function of two variables:
\begin{enumerate}
    \item \textbf{Likelihood (Probability):} Evaluated based on historical data, threat intelligence, and the asset's attractiveness to attackers.
    \item \textbf{Impact (Consequence):} Measured across the \textbf{CIA triad} (Confidentiality, Integrity, Availability). Prof. Atzeni emphasizes that impact is weighted by the criticality of the primary asset.
\end{enumerate}

\subsubsection*{3. The Risk Matrix and Treatment Strategies}
The intersection of impact and likelihood determines the risk level, which is compared against the \textbf{Risk Appetite}.
\begin{center}
\small
\begin{tabular}{|c|c|c|c|}
\hline
\textbf{Likelihood / Impact} & \textbf{Low} & \textbf{Medium} & \textbf{High} \\
\hline
\textbf{High} & Medium & High & \textbf{Critical} \\
\hline
\textbf{Medium} & Low & Medium & High \\
\hline
\textbf{Low} & \textit{Acceptable} & Low & Medium \\
\hline
\end{tabular}
\end{center}
Once risks are prioritized, four treatment options are available:
\begin{itemize}
    \item \textbf{Mitigation (Treatment):} Applying controls (e.g., encryption, MFA) to reduce risk to an acceptable level.
    \item \textbf{Avoidance:} Eliminating the risk by changing the business process.
    \item \textbf{Transfer (Sharing):} Moving the risk to a third party (e.g., insurance).
    \item \textbf{Acceptance:} Consciously deciding to keep the risk (management approval required).
\end{itemize}

\subsubsection*{4. Documentation: SoA and Residual Risk}
The process must be formal and auditable:
\begin{itemize}
    \item \textbf{Statement of Applicability (SoA):} A document identifying which ISO 27001 controls are applicable and how they are implemented.
    \item \textbf{Residual Risk:} The risk remaining after treatment. Prof. Atzeni notes that "zero risk" does not exist; residual risk must be formally accepted by the Board.
\end{itemize}

